\section{Theory}\label{sec:theory}
We formulate \COtwo solubility as a reacting electrolyte vapor--liquid-equilibrium problem.  At specified temperature, apparent amine concentration, and absorbed \COtwo loading, the calculation determines a charge-balanced true-species liquid composition and the corresponding vapor-phase \COtwo partial pressure.  The liquid composition enters the ePC-SAFT fugacity coefficients, activity coefficients, and electrostatic contributions.

\subsection{PC-SAFT residual Helmholtz energy}
ePC-SAFT is formulated as a residual Helmholtz-energy equation of state.  For a state defined by temperature \(T\), molar density \(\rho\), and mole-fraction vector \(\boldsymbol{x}\), the reduced residual Helmholtz energy \(\tilde{a}^{\mathrm{res}}=a^{\mathrm{res}}/(RT)\) is assembled from hard-chain, dispersion, association, Debye--Huckel, and Born terms \cite{Gross2001,Gross2002,Cameretti2005,Bulow2021a,Figiel2025}:
\begin{equation}
    \tilde{a}^{\mathrm{res}}
    =
    \tilde{a}^{\mathrm{hc}}
    +\tilde{a}^{\mathrm{disp}}
    +\tilde{a}^{\mathrm{assoc}}
    +\tilde{a}^{\mathrm{DH}}
    +\tilde{a}^{\mathrm{Born}} ,
    \label{eq:epcsaft-residual-split}
\end{equation}
The hard-chain contribution represents the repulsive chain reference, dispersion represents short-range attraction, association represents site-specific hydrogen bonding, the Debye--Huckel term represents long-range electrostatics, and the Born term represents electrostatic solvation.  In loaded \MEA, neutral solvent structure, ionic strength, and solvation all change with \COtwo loading.  Neutral \MEA and \HtwoO require segment, dispersion, and association parameters; charged species also require charge, dielectric, and Born-diameter information.

The hard-chain reference is built from packing-fraction moments
\begin{equation}
    \zeta_n=\frac{\pi}{6}\rho\sum_i x_i m_i d_i^n,\qquad n\in\{0,1,2,3\},
    \label{eq:epcsaft-zeta-moments}
\end{equation}
where \(m_i\) is the segment number and \(d_i\) is the temperature-dependent segment diameter.  With \(\bar{m}=\sum_i x_i m_i\), the chain correction subtracts the connectivity penalty from the hard-sphere reference,
\begin{equation}
    \tilde{a}^{\mathrm{hc}}
    =
    \bar{m}\tilde{a}^{\mathrm{hs}}
    -
    \sum_i x_i(m_i-1)\ln g_{ii}^{\mathrm{hs}}(\sigma_{ii}) .
    \label{eq:epcsaft-hard-chain}
\end{equation}
The hard-chain term sets the excluded-volume and density scales that enter fugacity and liquid-state activity coefficients.

The dispersion term introduces attractive perturbations through mixture moments of segment energy and size,
\begin{equation}
    \tilde{a}^{\mathrm{disp}}
    =
    -2\pi\rho I_1(\eta,\bar{m})\overline{m^2\epsilon\sigma^3}
    -\pi\rho\bar{m}C_1 I_2(\eta,\bar{m})
    \overline{m^2\epsilon^2\sigma^3},
    \label{eq:epcsaft-dispersion}
\end{equation}
Here \(\eta=\zeta_3\), \(I_1\) and \(I_2\) are the PC-SAFT dispersion integrals, and the overbars denote pair-mixed composition averages.  Binary interaction parameters modify the mixed dispersion energy \(\epsilon_{ij}\), which links pressure residuals to neutral--neutral and ion--molecule dispersion choices.

Association is represented through site fractions \(X^{A_i}\),
\begin{equation}
    \tilde{a}^{\mathrm{assoc}}
    =
    \sum_i x_i \sum_{A_i}
    \left[
        \ln X^{A_i}
        -\frac{X^{A_i}}{2}
        +\frac{1}{2}
    \right],
    \label{eq:epcsaft-association}
\end{equation}
with association strengths proportional to
\begin{equation}
    \Delta^{A_iB_j}_{ij}
    =
    d_{ij}^{3}g_{ij}^{\mathrm{hs}}
    \kappa^{A_iB_j}_{ij}
    \left[
        \exp\left(\frac{\epsilon^{A_iB_j}_{ij}}{k_{\mathrm{B}}T}\right)-1
    \right].
    \label{eq:epcsaft-association-strength}
\end{equation}
This contribution represents hydrogen bonding in water and alkanolamines.  Ions without assigned association sites contribute no association-site balances.

\subsection{Electrolyte PC-SAFT and ion activity coefficients}
For charged species, the Debye--Huckel contribution uses the dielectric state and the inverse Debye length,
\begin{align}
    \kappa
    &=
    \left[
        \frac{\rho e^2}{k_{\mathrm{B}}T\varepsilon_0\varepsilon_r}
        \sum_j x_j z_j^2
    \right]^{1/2},
    \label{eq:epcsaft-debye-kappa} \\
    \tilde{a}^{\mathrm{DH}}
    &=
    -\frac{\kappa e^2}{12\pi\varepsilon_0\varepsilon_r k_{\mathrm{B}}T}
    \sum_i x_i z_i^2\chi_i ,
    \label{eq:epcsaft-debye-helmholtz}
\end{align}
Here \(z_i\) is the ionic charge and \(\chi_i\) is the finite-ion-size correction.  Relative permittivity \(\varepsilon_r\) enters both the Debye--Huckel and Born terms.  Their residual chemical potentials determine the ePC-SAFT activity coefficients used in the liquid reaction equations.

\subsection{Modified Born term with solvation shell and dielectric saturation}
The Born contribution is evaluated in modular form:
\begin{equation}
    \tilde{a}^{\mathrm{Born}}
    =
    -\frac{e^2}{4\pi\varepsilon_0 k_{\mathrm{B}}T}
    \sum_i x_i z_i^2
    \sum_{\delta\in\mathcal{D}}
    \left(1-\frac{1}{\varepsilon_{r,m_\delta}}\right)
    D_{i,\delta}^{(m_\delta)} .
    \label{eq:epcsaft-born-modular}
\end{equation}
The active set \(\mathcal{D}\) contains the base Born term and may include solvation-shell-modeling (SSM) and dielectric-saturation (DS) corrections,
\begin{equation}
    \mathcal{D}
    =
    \{\mathrm{Born}\}\cup\mathcal{D}_{\mathrm{add}},
    \qquad
    \mathcal{D}_{\mathrm{add}}\subseteq\{\mathrm{SSM},\mathrm{DS}\}.
    \label{eq:epcsaft-born-active-set}
\end{equation}
For SSM+DS, \(\mathcal{D}_{\mathrm{add}}=\{\mathrm{SSM},\mathrm{DS}\}\).  The base Born and SSM terms use the bulk dielectric medium, while the DS term uses the ion-local dielectric medium:
\begin{equation}
    m_\delta=
    \begin{cases}
        \mathrm{bulk}, & \delta\in\{\mathrm{Born},\mathrm{SSM}\},\\
        \mathrm{ion}, & \delta=\mathrm{DS}.
    \end{cases}
    \label{eq:epcsaft-born-medium}
\end{equation}
The inverse-diameter terms are
\begin{align}
    D_{i,\mathrm{Born}}^{(\mathrm{bulk})}
    &=
    \frac{1}{d_i^{\mathrm{Born}}},
    \label{eq:epcsaft-born-base-dterm}\\
    D_{i,\mathrm{SSM}}^{(\mathrm{bulk})}
    &=
    \frac{1}{d_i^{\mathrm{Born}}+\Delta d_i}
    -
    \frac{1}{d_i^{\mathrm{Born}}},
    \label{eq:epcsaft-born-ssm-dterm}\\
    D_{i,\mathrm{DS}}^{(\mathrm{ion})}
    &=
    \frac{1}{d_i^{\mathrm{Born}}}
    -
    \frac{1}{d_i^{\mathrm{Born}}+\Delta d_i},
    \label{eq:epcsaft-born-ds-dterm}
\end{align}
with shell shift
\begin{equation}
    \Delta d_i
    =
    \frac{f_{\mathrm{mix}}-1}{|z_i|}\,d_i^{\mathrm{Born}},
    \qquad
    f_{\mathrm{mix}}=\sum_k x_k f_k .
    \label{eq:epcsaft-born-shell-shift}
\end{equation}
The SSM term changes the effective Born radius in the bulk medium, and the DS term assigns the complementary radius increment to the ion-local dielectric medium.  Carbonation changes both the ionic population and solvent environment, so both contributions vary with loading.

We use concentration-dependent bulk permittivity and the SSM+DS convention of Figiel et al.\ \cite{Figiel2025}.  That convention assigns pure ions \(\varepsilon_r=8\) and \(f_k=1\), while solvent \(f_k\) values are fitted to electrolyte activity data; reported values are 1.5 for water, 1.4 for methanol, and 1.6 for ethanol.  Hajj et al.\ measured loaded 30 wt\% \MEA over 915--2450 MHz and obtained a zero-frequency parameter by fitting a Cole--Cole dispersion model\cite{Hajj2024}.  The microwave extrapolation does not independently identify the \MEA solvent-shell parameter, so the present calculation keeps \(f_{\MEA}=1\).  With concentration-dependent permittivity, ion solvation responds to the changes in solvent composition and ionic strength produced by carbonation.  \Cref{tab:full-ionic-ssm-ds-parameters} lists the adopted \MEA-family inputs and electrolyte-literature conventions.

\subsection{MEA--CO2--H2O reaction network}
The unloaded \MEA mass fraction \(w_{\MEA}\) and \COtwo loading \(\alpha\), defined as moles of absorbed \COtwo per mole of total amine, specify the apparent feed.  Molecular-weight conversion gives the unloaded liquid composition:
\begin{equation}
    x_{\MEA,0}
    =
    \frac{w_{\MEA}}
    {M_{\MEA}/M_{\HtwoO}+w_{\MEA}\left(1-M_{\MEA}/M_{\HtwoO}\right)},
    \qquad
    x_{\HtwoO,0}=1-x_{\MEA,0}.
    \label{eq:unloaded-mea-composition}
\end{equation}
We add absorbed \COtwo according to \(n_{\COtwo,0}=\alpha n_{\MEA,0}\) and normalize the apparent state before solving the reactions.

The true-species liquid basis resolves the apparent \MEA and \COtwo inventories into the molecular and ionic species used for acid--base and solubility calculations:
\begin{equation}
    \boldsymbol{x}
    =
    \left(
    x_{\COtwo},x_{\MEA},x_{\HtwoO},x_{\MEAH},x_{\MEACOO},
    x_{\HCOthree},x_{\COthree},x_{\HthreeO},x_{\Hydroxide}
    \right).
    \label{eq:true-species-vector}
\end{equation}
The speciation calculation uses the five independent liquid reactions in \cref{tab:reaction-equilibrium-constants}.  For reaction \(r\), thermodynamic equilibrium is written in terms of activities,
\begin{equation}
    \sum_i \nu_{ri}\ln a_i = \ln K_r(T),
    \qquad
    a_i = \gamma_i^{\mathrm{ePC-SAFT}} x_i,
    \label{eq:activity-reaction-equilibrium}
\end{equation}
where $\nu_{ri}$ is the stoichiometric coefficient and $\gamma_i^{\mathrm{ePC-SAFT}}$ is the activity coefficient evaluated from the ePC-SAFT liquid state.  The reaction constants are represented as
\begin{equation}
    \ln K_r(T)=A_r+\frac{B_r}{T}+C_r\ln T+D_rT,
    \label{eq:reaction-equilibrium-constant-temperature}
\end{equation}
with \(T\) in kelvin and a mole-fraction activity convention.  \Cref{tab:reaction-equilibrium-constants} gives the source range of each coefficient correlation.  The ideal baseline replaces each activity coefficient with unity.

\begin{table*}[t]
    \centering
    \scriptsize\rmfamily
    \setlength{\tabcolsep}{3pt}
    \renewcommand{\arraystretch}{1.12}
    \caption{Liquid-phase reaction network and equilibrium-constant coefficients used for the activity-coupled speciation solve.  The constants use \(\ln K_r(T)=A_r+B_r/T+C_r\ln T+D_rT\), with \(T\) in kelvin and \(a_i=\gamma_i x_i\) on a mole-fraction activity basis.  Coefficients are taken from the acid-gas/ethanolamine reaction basis used in PC-SAFT calculations~\cite{Nasrifar2010}.}
    \label{tab:reaction-equilibrium-constants}
    \begin{adjustbox}{max width=\textwidth}
    \begin{tabularx}{\textwidth}{@{}c>{\raggedright\arraybackslash}Xcccccc@{}}
        \toprule
        ID & Liquid reaction & Source constant & \(A_r\) & \(B_r\) & \(C_r\) & \(D_r\) & Source range (K) \\
        \midrule
        R1 & \(2\HtwoO \rightleftharpoons \HthreeO+\Hydroxide\) & \(k_1^x\) & 132.899 & -13445.9 & -22.4773 & 0 & 273--498 \\
        R2 & \(\COtwo+2\HtwoO \rightleftharpoons \HCOthree+\HthreeO\) & \(k_2^x\) & 231.456 & -12092.1 & -36.7816 & 0 & 273--498 \\
        R3 & \(\HCOthree+\HtwoO \rightleftharpoons \COthree+\HthreeO\) & \(k_3^x\) & 216.049 & -12431.7 & -35.4819 & 0 & 273--498 \\
        R4 & \(\MEACOO+\HtwoO \rightleftharpoons \MEA+\HCOthree\) & \(k_9^x\) & 2.8898 & -3635.09 & 0 & 0 & 298--393 \\
        R5 & \(\MEAH+\HtwoO \rightleftharpoons \MEA+\HthreeO\) & \(k_7^x\) & 2.1211 & -8189.38 & 0 & -0.007484 & 273--323 \\
        \bottomrule
    \end{tabularx}
    \end{adjustbox}
\end{table*}


The liquid composition also satisfies carbon, amine, and water pseudo-component balances plus electroneutrality:
\begin{align}
    x_{\COtwo,0} &= x_{\COtwo}+x_{\MEACOO}+x_{\HCOthree}+x_{\COthree}, \\
    x_{\MEA,0} &= x_{\MEA}+x_{\MEAH}+x_{\MEACOO}, \\
    x_{\HtwoO,0} &= x_{\HtwoO}+x_{\HCOthree}+x_{\COthree}+x_{\HthreeO}+x_{\Hydroxide}, \\
    x_{\MEAH}+x_{\HthreeO} &= x_{\MEACOO}+x_{\HCOthree}+2x_{\COthree}+x_{\Hydroxide}.
    \label{eq:true-species-balances}
\end{align}

\subsection{Phase-equilibrium conditions and reference states}
For a liquid state at temperature \(T\), pressure \(P\), and mole-fraction vector \(\boldsymbol{x}\), ePC-SAFT provides residual chemical potentials and fugacity coefficients.  Molecular vapor--liquid equilibrium requires equal fugacities,
\begin{equation}
    f_i^{\liq}(T,P,\boldsymbol{x}^{\liq})
    =
    f_i^{\vap}(T,P,\boldsymbol{y}^{\vap}),
    \qquad i \in \{\COtwo,\HtwoO,\MEA\},
    \label{eq:fugacity-equality}
\end{equation}
while ions remain in the liquid phase.  Residual chemical potentials connect phase equilibrium to the activity coefficients in the reaction equations:
\begin{equation}
    \ln\varphi_k
    =
    \tilde{\mu}_k^{\mathrm{res}}-\ln Z,
    \qquad
    \ln\gamma_i^{*}
    =
    \ln\varphi_i-\ln\varphi_i^{\infty},
    \label{eq:epcsaft-phi-gamma-link}
\end{equation}
where \(\varphi_i^{\infty}\) is the infinite-dilution reference fugacity coefficient for the liquid activity coefficient.  This convention places neutral and ionic activities on the same true-species reaction basis.  Each residual Helmholtz contribution therefore affects both phase equilibrium and reaction equilibrium.

